%% kapa: exam=1 round=28 year=2017 period=2 form=A subject=appraisal_law subject_ko=감정평가관계법규 seq=36 number=36
%% source: Q-Net 감정평가사 기출문제 (cst00309) artlSeq=5199603 file=extracted/1008/5199603/2017년 제28회 감정평가사 1차시험 문제지/2교시 A형.pdf
%% conversion: kordoc 4.24.51 -> tools/gen1.py (auto); review=proofread-sonnet; flags=none
\begin{question}{36}
동산ㆍ채권 등의 담보에 관한 법령상 채권담보권에 관한 설명으로 옳지 않은 것은?
\begin{choices}
  \item 법인 등이 담보약정에 따라 금전의 지급을 목적으로 하는 지명채권을 담보로 제공하는 경우에는 담보등기를 할 수 있다.
  \item 채무자가 특정되지 아니한 여러 개의 채권이더라도 채권의 종류, 발생 원인, 발생 연월일을 정하는 등의 방법으로 특정할 수 있는 경우에는 이를 목적으로 하여 담보등기를 할 수 있다.
  \item 채권담보권의 목적이 된 채권이 피담보채권보다 먼저 변제기에 이른 경우에는 담보권자는 제3채무자에게 그 변제금액의 공탁을 청구할 수 있다.
  \item 담보권자는「민사집행법」에서 정한 집행방법으로는 채권담보권을 실행할 수 없다.
  \item 담보권자는 피담보채권의 한도에서 채권담보권의 목적이 된 채권을 직접 청구할 수 있다.
\end{choices}
\end{question}
\answer{4}
